\begin{abstract}
Pretrained image-conditioned 3D generators offer a practical starting point for
reconstructing objects from visual observations, but sparse patent drawings
present a substantially different input domain from natural photographs.
Thin technical contours, dashed boundaries, annotations, and multiple depicted
views provide limited appearance information and leave depth and unseen
geometry ambiguous. In this work, we investigate input-level adaptation of a
frozen TRELLIS generator to these sparse, black-and-white sketches. Our
implemented sketch-aware stage constructs deterministic distance-to-ink maps
that expand the spatial support of technical linework, with optional Gaussian
smoothing to control the transition around detected strokes. This formulation
introduces sketch-derived structural cues without retraining the pretrained
backbone. We implement four matched configurations that separate standard
image preprocessing, bypassed background removal, distance mapping, and
smoothing, using shared reviewed crops and paired generation seeds.
To characterize the reconstruction setting, we use patent TIFF drawings from
IMPACT to generate 438 three-dimensional reconstructions associated with
50 patents using pretrained TRELLIS. We assess these generated assets through
multi-view image comparisons using structural similarity, peak
signal-to-noise ratio, learned perceptual distance, and CLIP image and caption
similarity, alongside recomputed mesh statistics and qualitative inspection.
Scores are aggregated across rendered views, within drawings, and then within
patents to avoid giving greater weight to patents with more drawings. A
view-coverage diagnostic examines how the evaluated camera angles affect
reported similarity, distinguishing favorable projection selection from
reconstruction quality. Inspection of saved input--render pairs further
identifies limitations of threshold-based foreground extraction for thin
linework. The completed measurements establish a descriptive baseline for
our patent reconstructions and motivate a controlled evaluation of input
adaptation; comparative adaptation gains remain to be established through
the matched generation experiments.
\end{abstract}
